\begin{thebibliography}{99}
\bibitem{Arratia98} R. Arratia - “On the central role of scale invariant Poisson processes on (0,∞)”, in Microsurveys in discrete probability (Princeton, NJ, 1997), DIMACS Ser. Discrete Math. Theoret. Comput. Sci., vol. 41, American Mathematical Society, Providence, RI, 1998, p. 21-41
\bibitem{Arratia02} R. Arratia - “On the amount of dependence in the prime factorization of a uniform random integer”, in Contemporary combinatorics, Bolyai Soc. Math. Stud., vol. 10, János Bolyai Math. Soc., Budapest, 2002, p. 29-91
\bibitem{ABT06} R. Arratia, A. D. Barbour \& S. Tavaré - “A tale of three couplings: Poisson-Dirichlet and GEM approximations for random permutations”, Combin. Probab. Comput. 15 (2006) no. 1-2, p. 31-62
\bibitem{BareikisManstavicius07} G. Bareikis \& E. Manstavičius - “On the DDT theorem”, Acta Arith. 126 (2007) no. 2, p. 155-168
\bibitem{Bill72} P. Billingsley - “On the distribution of large prime divisors”, Period. Math. Hungar. 2 (1972), p. 283-289
\bibitem{delaBretecheTenenbaum16} R. de la Bretèche \& G. Tenenbaum - “Sur les processus arithmétiques liés aux diviseurs”, Adv. in Appl. Probab. 48 (2016) no. A, p. 63-76
\bibitem{DDT79} J.-M. Deshouillers, F. Dress \& G. Tenenbaum - “Lois de répartition des diviseurs. I”, Acta Arith. 34 (1979) no. 4, p. 273-285 (loose errata)
\bibitem{DonnellyTavare87} P. Donnelly \& S. Tavaré - “The population genealogy of the infinitely-many neutral alleles model”, J. Math. Biol. 25 (1987) no. 4, p. 381-391
\bibitem{Feng10} S. Feng - The Poisson-Dirichlet distribution and related topics. Models and asymptotic behaviors, Probability and its Appl., Springer, Heidelberg, 2010
\bibitem{JKB} N. L. Johnson, S. Kotz \& N. Balakrishnan - Discrete multivariate distributions, Wiley Series in Probab. and Stat.: Applied Probab. and Stat., John Wiley \& Sons, Inc., New York, 1997
\bibitem{Kingman93} J. F. C. Kingman - Poisson processes, Oxford Studies in Probab., vol. 3, The Clarendon Press, Oxford University Press, New York, 1993
\bibitem{KotzBal00} S. Kotz, N. Balakrishnan \& N. L. Johnson - Continuous multivariate distributions. Vol. 1. Models and applications, Wiley Series in Probab. and Stat.: Applied Probab. and Stat., Wiley-Interscience, New York, 2000
\bibitem{Koukoulo19} D. Koukoulopoulos - The distribution of prime numbers, Graduate Studies in Math., vol. 203, American Mathematical Society, Providence, RI, 2019
\bibitem{LastPenrose18} G. Last \& M. Penrose - Lectures on the Poisson process, Institute of Math. Stat. Textbooks, vol. 7, Cambridge University Press, Cambridge, 2018
\bibitem{Ken22} S.-K. Leung - “Dirichlet law for factorisation of integers, polynomials and permutations”, Math. Proc. Cambridge Philos. Soc. 175 (2023) no. 3, p. 649-676
\bibitem{KenWeb24} S.-K. Leung - “Dirichlet law with arbitrary parameters for factorizations of integers” (2025), Retrieved on 24 August 2025 from https://sites.google.com/view/sunkaileung/miscellaneous
\bibitem{MV07} H. L. Montgomery \& R. C. Vaughan - Multiplicative number theory. I. Classical theory, Cambridge Studies in Advanced Math., vol. 97, Cambridge University Press, Cambridge, 2007
\bibitem{NyandwiSmati13} S. Nyandwi \& A. Smati - “Distribution laws of pairs of divisors”, Integers 13 (2013), article ID A13, 13 pages
\bibitem{Tenenbaum00} G. Tenenbaum - “A rate estimate in Billingsley’s theorem for the size distribution of large prime factors”, Q. J. Math. 51 (2000) no. 3, p. 385-403
\end{thebibliography}
